\section{2. Directory Structure}\label{directory-structure}

This chapter is the canonical reference for the repository layout. Every
top-level folder has a single, well-defined purpose. Understanding this
structure is a prerequisite for all other work.

\subsection{2.1 Full Layout}\label{full-layout}

\begin{minted}{text}
experimentos_mestrado_imbalnced_ts/
│
├── README.md                     <- Master dashboard (auto-updated, never edit by hand)
├── CONTRIBUTING.md               <- The rules: how to add / maintain experiments
├── .gitignore                    <- What Git ignores: venvs, datasets, build artifacts
│
├── experiments/                  <- One subfolder per paper replication
│   │
│   ├── _template/                <- Blueprint. NEVER edit directly.
│   │   ├── README.md             <- Structured experiment report (fill this in)
│   │   ├── REPLICATION_LOG.md    <- Daily work diary (prepend new entries)
│   │   ├── requirements.txt      <- Python dependencies for this experiment
│   │   ├── config/
│   │   │   └── default.yaml      <- Key parameters: seeds, paths, hyperparams
│   │   ├── src/
│   │   │   ├── original/         <- AUTHORS' code goes here (manual, never scripted)
│   │   │   └── adapted/          <- YOUR changes and wrappers go here
│   │   ├── notebooks/            <- Exploratory Jupyter notebooks
│   │   └── results/
│   │       ├── tables/           <- CSV / XLSX output tables
│   │       └── figures/          <- Generated plots and visualizations
│   │
│   ├── 001_uba_utility_regression/
│   ├── 002_ts_resampling_strategies/
│   ├── ... (17 experiment folders in total)
│   └── 017_imbalanced_regression_pipeline/
│
├── scripts/                      <- Automation tools (bookkeeping only)
│   ├── new_experiment.sh         <- Scaffolds a new experiment from _template
│   ├── update_readme.py          <- Rebuilds the progress dashboard in README.md
│   └── sync_from_excel.py        <- Syncs paper metadata from Google Drive Excel
│
├── shared/                       <- Code reused across multiple experiments
│   ├── datasets/
│   │   └── README.md             <- Registry of shared datasets and download instructions
│   ├── metrics/
│   │   ├── regression.py         <- RMSE, MAE, SMAPE — fallback only; prefer paper's own
│   │   └── classification.py     <- F1, Precision, Recall — fallback only
│   ├── visualization/
│   │   └── plot_utils.py         <- Publication-quality matplotlib helpers
│   └── utils/
│       └── data_loading.py       <- Path helpers for navigating the repo structure
│
├── docs/                         <- Reference documents and this guide
│   ├── guide/                    <- Source for this PDF
│   ├── paper_registry.md         <- Auto-generated registry of all 43 papers
│   └── methodology.md            <- Scientific replication methodology
│
└── .github/
    └── ISSUE_TEMPLATE/
        └── new-experiment.md     <- GitHub issue template for proposing new experiments
\end{minted}

\subsection{2.2 Folder Reference}\label{folder-reference}

\subsubsection{\texorpdfstring{\texttt{experiments/} --- The Core of the
Repository}{experiments/ --- The Core of the Repository}}\label{experiments-the-core-of-the-repository}

This is where all research lives. Each subfolder corresponds to one
paper replication and follows a strict naming convention:

\begin{minted}{text}
NNN_descriptive_name/
\end{minted}

Where \texttt{NNN} is a zero-padded three-digit number
that encodes the \textbf{priority order} of the experiment. Lower
numbers are done first.

\textbf{Golden rules for \texttt{experiments/}:}

\begin{itemize}
\tightlist
\item
  Never add a folder by hand. Always use
  \texttt{new\textbackslash\{\}\_experiment.sh}.
\item
  Never rename an existing folder. The number is the experiment's
  permanent identifier.
\item
  The folder structure inside each experiment must match
  \texttt{\textbackslash\{\}\_template/} exactly.
\end{itemize}

\subsubsection{\texorpdfstring{\texttt{experiments/\_template/} --- The
Blueprint}{experiments/\_template/ --- The Blueprint}}\label{experiments_template-the-blueprint}

This folder defines the schema for every experiment in the repository.
The automation script \texttt{new\textbackslash\{\}\_experiment.sh}
copies it in its entirety when creating a new experiment.

\begin{quote}
\small \textbf{Critical rule:} Never edit \texttt{\textbackslash\{\}\_template/}
to fix a single experiment. If you edit the template, you are changing
the contract for all \emph{future} experiments. Edit the specific
experiment folder instead.
\end{quote}

\textbf{Contents of each experiment folder:}

{\def\LTcaptype{none} % do not increment counter
\small
\begin{longtable}[]{@{}
  >{\raggedright\arraybackslash}p{(\linewidth - 4\tabcolsep) * \real{0.3333}}
  >{\raggedright\arraybackslash}p{(\linewidth - 4\tabcolsep) * \real{0.3333}}
  >{\raggedright\arraybackslash}p{(\linewidth - 4\tabcolsep) * \real{0.3333}}@{}}
\toprule\noalign{}
\begin{minipage}[b]{\linewidth}\raggedright
File / Folder
\end{minipage} & \begin{minipage}[b]{\linewidth}\raggedright
Purpose
\end{minipage} & \begin{minipage}[b]{\linewidth}\raggedright
Who fills it
\end{minipage} \\
\midrule\noalign{}
\endhead
\bottomrule\noalign{}
\endlastfoot
\texttt{README.md} & Structured scientific report:
metadata, results, status & You \\
\texttt{REPLICATION\textbackslash\{\}\_LOG.md} & Chronological work
diary, newest entry on top & You \\
\texttt{requirements.txt} & Exact Python package
versions (\texttt{pip freeze}) & You \\
\texttt{config/default.yaml} & Experiment parameters:
seeds, paths, hyperparams & You \\
\texttt{src/original/} & Authors' original source code,
cloned or copied & You (manual) \\
\texttt{src/adapted/} & Your modifications, wrappers,
and analysis code & You \\
\texttt{notebooks/} & Jupyter notebooks for exploration
& You \\
\texttt{results/tables/} & Output CSVs and comparison
tables & Experiments \\
\texttt{results/figures/} & Generated plots &
Experiments \\
\end{longtable}
}

\subsubsection{\texorpdfstring{\texttt{scripts/} --- Automation
Tools}{scripts/ --- Automation Tools}}\label{scripts-automation-tools}

These scripts handle \textbf{bookkeeping only}. They never touch your
data or experiment results.

{\def\LTcaptype{none} % do not increment counter
\small
\begin{longtable}[]{@{}
  >{\raggedright\arraybackslash}p{(\linewidth - 4\tabcolsep) * \real{0.3333}}
  >{\raggedright\arraybackslash}p{(\linewidth - 4\tabcolsep) * \real{0.3333}}
  >{\raggedright\arraybackslash}p{(\linewidth - 4\tabcolsep) * \real{0.3333}}@{}}
\toprule\noalign{}
\begin{minipage}[b]{\linewidth}\raggedright
Script
\end{minipage} & \begin{minipage}[b]{\linewidth}\raggedright
When to run
\end{minipage} & \begin{minipage}[b]{\linewidth}\raggedright
What it does
\end{minipage} \\
\midrule\noalign{}
\endhead
\bottomrule\noalign{}
\endlastfoot
\texttt{new\textbackslash\{\}\_experiment.sh} & Adding an experiment
after \#017 & Copies \texttt{\textbackslash\{\}\_template/}, renames
placeholders, sets date \\
\texttt{update\textbackslash\{\}\_readme.py} & After changing any
experiment's status & Rebuilds the status table in the root
\texttt{README.md} \\
\texttt{sync\textbackslash\{\}\_from\textbackslash\{\}\_excel.py} & After updating the
Google Drive paper spreadsheet & Regenerates
\texttt{docs/paper\textbackslash\{\}\_registry.md} \\
\end{longtable}
}

\subsubsection{\texorpdfstring{\texttt{shared/} --- Cross-Experiment
Code}{shared/ --- Cross-Experiment Code}}\label{shared-cross-experiment-code}

Code that is genuinely reused across multiple experiments lives here
instead of being duplicated inside each experiment folder.

\textbf{Important usage rule:} The
\texttt{shared/metrics/} implementations are
\textbf{fallback implementations} only. If the paper provides its own
metric code, use that. Shared metrics exist to avoid duplication when no
reference implementation is available.

\subsubsection{\texorpdfstring{\texttt{docs/} ---
Documentation}{docs/ --- Documentation}}\label{docs-documentation}

The \texttt{docs/} folder contains reference material
about the project as a whole, not about any individual experiment.

{\def\LTcaptype{none} % do not increment counter
\small
\begin{longtable}[]{@{}
  >{\raggedright\arraybackslash}p{(\linewidth - 2\tabcolsep) * \real{0.5000}}
  >{\raggedright\arraybackslash}p{(\linewidth - 2\tabcolsep) * \real{0.5000}}@{}}
\toprule\noalign{}
\begin{minipage}[b]{\linewidth}\raggedright
File
\end{minipage} & \begin{minipage}[b]{\linewidth}\raggedright
Description
\end{minipage} \\
\midrule\noalign{}
\endhead
\bottomrule\noalign{}
\endlastfoot
\texttt{guide/} & Source Markdown files for this PDF
guide \\
\texttt{paper\textbackslash\{\}\_registry.md} & Auto-generated list of
all 43 papers from the Excel spreadsheet \\
\texttt{methodology.md} & The scientific replication
methodology (what counts as a valid replication) \\
\end{longtable}
}

\subsection{2.3 The Two Critical
Boundaries}\label{the-two-critical-boundaries}

These two rules are the most important structural decisions in the
repository. Violating them breaks traceability.

\textbf{Boundary 1 --- \texttt{\textbackslash\{\}\_template/} is
immutable per experiment.} The script creates your experiment folder
from the template. After that, only \emph{you} edit the experiment
folder. The template is never touched again for that experiment.

\textbf{Boundary 2 --- \texttt{src/original/} is always
filled manually.} No script will ever write to
\texttt{src/original/}. You clone or copy the authors'
code there. This makes the provenance of all code 100\% explicit and
auditable.
